% Pursuing Stacks (cotes 134-2 to 134-8): archive cote 134-5, site page 65 (PDF page 66), typescript page 252.
% typescript. Draft compared: ps4.tex:3264-3339.
% AI-transcribed by Claude Opus 5.5 (claude-opus-5-5, adaptive thinking, effort high), Message Batches API, batch ps2 (msgbatch_01NVwr1Yuo8qJf5rykeGBXPK), 2026-11-12; not yet checked by eye.
\scanpage{134-5/65}
\mspage{252}

H 4)\quad u transforms any contractible object\inserted{x} of M into a contractible
object of M' .

H 4')\quad Same as H 4) , with x being restricted to be in a given \xout{\illegible{1}}
subcategory C of M, generating for the contractibility struc-%
ture of M.

\uline{Remark}\hand{ 2.} It should be noted that the condition\inserted{upon C} stated\deleted{,} in H 4') \deleted{upon C,} namel\hand{y}
that the (given) C should be generating for the contractibility structu\xout{r}%
re $M_c$ of M, means exactly two things (a) $C \hand{\subset} M_c$ , and (b) the fami-%
ly of all intervals $\underline{I} = (I,\hand{\delta}_0,\hand{\delta}_1)$ made up with objects of C (these in-%
tervals are necessarily homotopy intervals for the\xout{\illegible{1}} homotopy struc-%
ture of M) \uline{generates} the homotopy structure of M, i.e. two arrows in M
are homotopic iff they can be joined by a chain a arrows, two consecuti-%
ve among which being related by an elementary homotopy involving an
interval of that family. This \xout{\illegible{1}}\inserted[typed]{reminder} being made, it follows that C
\xout{can} is equally eligible for applying criterion H 3'),\unsure{,}which means that
\xout{in H 4')} we get still another equivalent formulation of H 4'), by deman-%
ding merely that for any two sections $\hand{\delta}_0,\hand{\delta}_1$ of x over $e_M$ , the corres-%
ponding sections of u(x) should be homotopic.

After these preliminaries, we can state at last \xout{what seems to me}
the following generalisation of the main result of section 65 (th.1
p. 176), concerning test functors with values in (Cat):

\margin{$\times$ \unsure{small unnecessary}, \unsure{cf. also remarks} on p. \unsure{348}, \illegible{1}}
\uline{Theorem 1} . Let $(M,M_c)$ be a\xout{n} contractibility structure\deleted{,}\inserted{(cf p. 248)}\hand{,} C a \unsure{small} full
\margin{\illegible{3}}
subcategory of M generating the contractibility structure (cf. remark 2
above), A a small category, and i : $A \longrightarrow M$ a functor, factoring
through $M_c$ . We consider M as endowed equally with the asphericity
structure $M_{as}$ generated by $M_c$ (cf. prop.2). Then the following condi-%
tions on i are equivalent:

(i) The functor $i^{*}$: $M \to A^{\wedge}$ deduced from i is compatible with the
homotopy structures on M and on $A^{\wedge}$ (cf. page 250\inserted{-251} for the latter), i.e.
$i^{*}$ satisfies ei\hand{$\downarrow$}ther one of the six equivalent conditions H 1) to H 4')
above.

(ii) The functor $i^{*}$ is locally $M_{as}$-aspheric (def.2 p.245), i.e.
for any x in $M_{as}$ , $i^{*}(x)$ is a locally aspheric object of $A^{\wedge}$ (cf. p.250),
i.e. $i^{*}(x)$ is aspheric over $e_{A^{\wedge}}$ .

(iii) \xout{For any}\inserted[typed]{Same as} (ii), but with x restricted to be in C.

\uline{Corollary 1} . In order for i to be totally $M_{as}$-aspheric (def.2 p.245)

\margin{65}
